% ch03-auth
\section{التوثيق والصلاحيّات}

في الفصلين السابقين عرفنا الطبقات و\en{CQRS} والمنافذ (\en{ports}). في هذا الفصل
ننتقل إلى أول ما يقابله أي مستخدم عند فتح \en{ScholarPath}: كيف يُثبت أنه \emph{هو}،
وكيف يقرّر النظام بعد ذلك ما \emph{يُسمح} له بعمله. الموضوعان متلازمان لكنهما مختلفان
تمامًا، ومن المهم التفريق بينهما من البداية.

\begin{idea}[الفكرة: توثيق مقابل تفويض]
\textbf{التوثيق (\en{Authentication})} يجيب على سؤال: «مَن أنت؟» — أي التأكّد من
هويّة المستخدم عبر كلمة مرور أو حساب \en{Google} أو رمز.
أمّا \textbf{التفويض (\en{Authorization})} فيجيب على سؤال آخر: «طيّب، وأنت مين، تقدر
تعمل إيه؟» — أي هل يحقّ لهذا المستخدم رؤية هذا المورد أو تعديله؟
النظام ينفّذ الأول عند الدخول (ويصدر رمزين)، ثم ينفّذ الثاني في \emph{كل} طلب لاحق.
\end{idea}

كل ما سنشرحه هنا مأخوذ من الكود الحيّ داخل \code{server/src}، تحديدًا من طبقة
\en{Infrastructure} (الخدمات) وطبقة \en{Application} (المُعالِجات / \en{handlers}).

% ====================================================================
\subsection{الهويّة عبر ASP.NET Core Identity}

بدلًا من كتابة نظام مستخدمين من الصفر (وهو أمر خطير أمنيًّا)، يبني \en{ScholarPath}
على \en{ASP.NET Core Identity}: مكتبة مايكروسوفت الناضجة لإدارة الحسابات وكلمات المرور
والأدوار. نقطة الوصل هي الصنف \code{ApplicationUser} الذي \emph{يرث} من
\code{IdentityUser<Guid>}، فيأخذ كل ما توفّره \en{Identity} (كلمة المرور المُجزّأة،
عدّاد المحاولات الفاشلة، تأكيد البريد) ويضيف فوقه حقول العمل الخاصّة بنا.

\begin{note}
اخترنا \code{Guid} كنوع للمُعرِّف بدل الرقم المتسلسل: المُعرِّفات العشوائية لا تُخمَّن
ولا تكشف عدد المستخدمين، وتناسب التوزيع دون تصادم.
\end{note}

هنشوف تعريف \code{ApplicationUser} (مختصرًا) من ملف \code{Identity.cs}:

\begin{english}
\begin{lstlisting}[language={[Sharp]C}]
public class ApplicationUser : IdentityUser<Guid>, ISoftDeletable
{
    public string FirstName { get; set; } = default!;
    public string LastName  { get; set; } = default!;
    public AccountStatus AccountStatus { get; set; } = AccountStatus.Unassigned;
    public bool IsOnboardingComplete { get; set; }
    public DateTimeOffset? LastLoginAt { get; set; }

    // Dual-role support: the role active for the current session
    public string? ActiveRole { get; set; }

    public ICollection<RefreshToken> RefreshTokens { get; } = [];
    public ICollection<LoginAttempt> LoginAttempts { get; } = [];

    public string FullName => $"{FirstName} {LastName}".Trim();
}
\end{lstlisting}
\end{english}

نقرأ الأسطر المهمّة واحدة واحدة:
\begin{itemize}
\item \code{: IdentityUser<Guid>} — الوراثة التي تجلب البنية التحتية للهوية جاهزة،
  فلا نُخزّن كلمة المرور بأنفسنا أبدًا.
\item \code{AccountStatus} و\code{IsOnboardingComplete} — حقول عمل تُقرّر مسار
  المستخدم بعد الدخول (مثلًا: هل أكمل خطوات التسجيل؟).
\item \code{ActiveRole} — لأن المستخدم قد يملك أكثر من دور، نحفظ الدور الفعّال للجلسة.
\item \code{RefreshTokens} — مجموعة رموز التحديث المرتبطة بالمستخدم؛ سنعود إليها بعد قليل.
\end{itemize}

\subsubsection{تجزئة كلمة المرور}
النظام لا يخزّن كلمة المرور نصًّا صريحًا، بل \emph{تجزئتها} فقط عبر المنفذ
\code{IPasswordHasher}، وتنفيذه \code{IdentityPasswordHasher} يفوّض العمل إلى
\code{PasswordHasher} من \en{Identity} (خوارزمية \en{PBKDF2} مع مِلح
عشوائي لكل تجزئة). عند الدخول نستدعي \code{Verify} الذي يعيد النتيجة
دون أن يفكّ التجزئة أبدًا (التجزئة عملية باتجاه واحد).

\begin{secnote}
لماذا مِلح (\en{salt}) لكل مستخدم؟ لأنه يجعل تجزئتين لكلمتَي مرور متطابقتين
\emph{مختلفتين} في القاعدة، فيُبطل هجمات جداول \en{rainbow} الجاهزة، ويمنع
معرفة أن مستخدِمَين يستعملان الكلمة نفسها. ولأن التجزئة خلف المنفذ
\code{IPasswordHasher}، يمكن استبدال الخوارزمية مستقبلًا دون لمس بقية الكود.
\end{secnote}

\subsubsection{القفل بعد المحاولات الفاشلة}
لصدّ هجمات التخمين، يعدّ \code{LoginCommandHandler} المحاولات الفاشلة
\emph{المتتالية} داخل نافذة زمنية؛ فإذا بلغت الحدّ يُقفَل الحساب مؤقّتًا. هنشوف
قلب هذا المنطق:

\begin{english}
\begin{lstlisting}[language={[Sharp]C}]
private const int MaxFailedAttempts = 5;
private static readonly TimeSpan LockoutDuration = TimeSpan.FromMinutes(30);
// ...
var priorFailures = await db.LoginAttempts
    .CountAsync(a => a.Email == normalizedEmail && !a.Succeeded
                     && a.OccurredAt >= windowStart, ct);

if (priorFailures + 1 >= MaxFailedAttempts)
{
    user.LockoutEnd = now + LockoutDuration;   // lock for 30 minutes
    user.AccessFailedCount = priorFailures + 1;
}
\end{lstlisting}
\end{english}

الشرح:
\begin{itemize}
\item \code{MaxFailedAttempts = 5} — بعد خمس محاولات فاشلة متتالية يُقفَل الحساب.
\item \code{CountAsync(...)} — يحسب الفشل \emph{غير الناجح} منذ بداية النافذة فقط،
  فأي دخول ناجح يصفّر العدّاد عمليًّا.
\item \code{user.LockoutEnd = now + LockoutDuration} — يضبط لحظة انتهاء القفل (30 دقيقة)؛
  وأي محاولة قبلها تُرفَض مباشرة برسالة «الحساب مقفول مؤقّتًا».
\end{itemize}

كل محاولة — ناجحة أو فاشلة — تُسجَّل كصفٍّ في \code{LoginAttempts} مع الوقت وعنوان
\en{IP}، فيصبح لدينا سِجلّ يمكن تدقيقه لاحقًا. أمّا \textbf{تأكيد البريد} فتديره
\en{Identity} عبر خاصية \code{EmailConfirmed}؛ فحساب لم يؤكّد بريده يظل محدود
الصلاحيات حتى يضغط رابط التفعيل.

% ====================================================================
\subsection{الرمزان: access JWT و refresh token}

بعد نجاح التوثيق يُصدر \code{TokenService} \emph{زوجًا} من الرموز، لكل منهما دور مختلف:

\begin{center}
\begin{tabular}{L{3.2cm} L{11.4cm}}
\toprule
\textbf{\color{spPrimary}الرمز} & \textbf{\color{spPrimary}دوره وعمره} \\
\midrule
\en{Access JWT} & رمز حامل (\en{bearer}) قصير العمر جدًّا — \textbf{15 دقيقة} —
يُرفَق مع كل طلب في ترويسة \en{Authorization}. يحمل هويّة المستخدم وأدواره موقّعةً. \\
\en{Refresh token} & رمز عشوائي أطول عمرًا (7 أيام، أو 30 مع «تذكُّرني») وظيفته
\emph{الوحيدة} استبدال رمز الوصول المنتهي دون إعادة إدخال كلمة المرور. \\
\bottomrule
\end{tabular}
\end{center}

\begin{why}
لماذا رمز وصول قصير جدًّا؟ لأن الـ\en{JWT} — بمجرد إصداره — لا يمكن سحبه فعليًّا حتى
ينتهي (فحصه يتم دون الرجوع للقاعدة). فلو سُرِّب، نريد نافذة خطر ضيّقة: 15 دقيقة فقط.
والثمن (تجديد متكرّر) رخيص لأن التجديد يتم بصمت عبر رمز التحديث دون إزعاج المستخدم.
\end{why}

\subsubsection{كيف يُولَّد رمز التحديث ويُخزَّن؟}
النقطة الحسّاسة: رمز التحديث \emph{لا} يُخزَّن كما هو في القاعدة، بل تُخزَّن
\emph{تجزئته} فقط في الحقل \code{RefreshToken.TokenHash}. هنشوف التوليد والتجزئة من
\code{TokenService}:

\begin{english}
\begin{lstlisting}[language={[Sharp]C}]
private static string GenerateSecureRefreshToken()
{
    var bytes = RandomNumberGenerator.GetBytes(64);   // 64 crypto-random bytes
    return Convert.ToBase64String(bytes);
}

private static string HashRefreshToken(string token)
{
    var bytes = SHA256.HashData(Encoding.UTF8.GetBytes(token));
    return Convert.ToHexString(bytes);                // store the hash, not the token
}
\end{lstlisting}
\end{english}

الشرح:
\begin{itemize}
\item \code{RandomNumberGenerator.GetBytes(64)} — مولّد عشوائي \emph{تشفيري} (لا
  \code{Random} العادي)، فالرمز غير قابل للتخمين.
\item الرمز الخام يُرسَل للعميل فقط. أمّا ما يُحفَظ في القاعدة فهو ناتج
  \code{HashRefreshToken} — تجزئة \en{SHA-256} بصيغة \en{hex}.
\end{itemize}

\begin{keyp}[نقطة مهمة: تسريب القاعدة لا يعطي رمزًا صالحًا]
لو تسرّبت قاعدة البيانات كاملة، لن يجد المهاجم إلا \emph{تجزئات} لا يمكن عكسها إلى
الرمز الأصلي. وبما أن النظام يقارن بتجزئة ما يُقدَّم إليه، فالتجزئة المسروقة عديمة
الفائدة. نفس المبدأ نطبّقه على رمز إعادة تعيين كلمة المرور (سيأتي بعد قليل).
\end{keyp}

\subsubsection{التدوير عند كل استخدام}
رمز التحديث \textbf{يُدوَّر}: كل مرة يُستخدَم فيها، يُبطَل القديم فورًا ويُصدَر بديل
جديد. وإذا عاد رمزٌ مُبطَل للظهور مرّة أخرى، فهذه إشارة سرقة، فنُلغي \emph{كل} جلسات
المستخدم احتياطًا. هنشوف \code{RotateRefreshTokenAsync} (مختصرًا):

\begin{english}
\begin{lstlisting}[language={[Sharp]C}]
// Reuse detection: a revoked token presented again means it was already
// rotated -- a replay attack or a stale client. Revoke the whole family
// so a stolen token can never be exchanged.
if (stored.IsRevoked)
{
    await RevokeAllForUserAsync(stored.UserId, "Refresh token reuse detected", ct);
    return null;
}

// Otherwise: mark the old one revoked and issue a fresh replacement pair.
stored.IsRevoked = true;
stored.RevokedAt = clock.UtcNow;
stored.RevokedReason = "Rotated";
var pair = IssueTokens(stored.User, roles, stored.User.ActiveRole, rememberMe);
stored.ReplacedByTokenId = newToken.Id;   // chain old -> new for traceability
\end{lstlisting}
\end{english}

الشرح:
\begin{itemize}
\item \code{if (stored.IsRevoked)} — لو الرمز المُقدَّم مُبطَل سلفًا فهذا «إعادة
  استخدام»؛ نستدعي \code{RevokeAllForUserAsync} فنقطع كل جلسات المستخدم.
\item \code{stored.IsRevoked = true} — نُبطِل القديم أوّلًا، ثم نُصدر زوجًا جديدًا،
  فلا يبقى رمزان صالحان في وقت واحد.
\item \code{ReplacedByTokenId} — نربط القديم بالجديد، فيصير لدينا سلسلة يمكن تتبّعها.
\end{itemize}

\subsubsection{نفس الحماية لرمز إعادة تعيين كلمة المرور}
حقل \code{PasswordResetToken.TokenHash} يتبع نفس الفلسفة: يُولَّد رمز خام، يُرسَل
بالبريد، ويُخزَّن \emph{مُجزَّأً} فقط، صالحًا لساعة واحدة ولمرّة واحدة:

\begin{english}
\begin{lstlisting}[language={[Sharp]C}]
var rawToken = Convert.ToHexString(RandomNumberGenerator.GetBytes(32));
db.PasswordResetTokens.Add(new PasswordResetToken
{
    UserId    = user.Id,
    TokenHash = Hash(rawToken),        // only the hash is persisted
    ExpiresAt = now + TokenLifetime,   // 1 hour, single use (UsedAt)
});
// the raw token travels only inside the emailed link, never to the DB
\end{lstlisting}
\end{english}

لاحظ أيضًا أن معالج «نسيت كلمة المرور» \emph{لا} يكشف إن كان البريد مسجّلًا أم لا —
يعيد نجاحًا صامتًا دائمًا — حتى لا يُستغَل لاكتشاف الحسابات الموجودة.

% ====================================================================
\subsection{خوارزميّة التوقيع: RS256 مقابل HS256}

الـ\en{JWT} يُوقَّع كي لا يستطيع أحد تزييف محتواه. أمامنا خياران شائعان، واختار
\en{ScholarPath} الأأمن منهما:

\begin{center}
\begin{tabular}{L{3.4cm} L{5.5cm} L{5.5cm}}
\toprule
\textbf{\color{spPrimary}الوجه} & \textbf{\color{spPrimary}\en{HS256} (متماثل)} & \textbf{\color{spPrimary}\en{RS256} (غير متماثل)} \\
\midrule
المفتاح & سرّ واحد مشترك يُوقّع ويتحقّق & زوج: مفتاح خاصّ يوقّع، وعامّ يتحقّق \\
مَن يحتاج السرّ؟ & كل طرف يتحقّق يحتاج السرّ نفسه & المتحقّق يحتاج المفتاح العامّ فقط \\
خطر التسريب & مَن يعرف السرّ يستطيع \emph{تزوير} الرموز & المفتاح العامّ يتحقّق ولا يوقّع \\
التدوير & صعب (السرّ منتشر في كل مكان) & سهل عبر منفذ مفاتيح مستقلّ \\
\bottomrule
\end{tabular}
\end{center}

في \code{TokenService} نوقّع بالمفتاح الخاصّ الآتي من منفذ المفاتيح
\code{IJwtKeyProvider}:

\begin{english}
\begin{lstlisting}[language={[Sharp]C}]
// RS256: sign with the RSA private key supplied by the key provider
// (Azure Key Vault in production, a local key in development).
var signingKey = keyProvider.GetSigningKey();
var creds = new SigningCredentials(signingKey, SecurityAlgorithms.RsaSha256);

var token = new JwtSecurityToken(
    issuer:   _opts.Issuer,
    audience: _opts.Audience,
    claims:   claims,
    notBefore: clock.UtcNow.UtcDateTime,
    expires:   expires.UtcDateTime,
    signingCredentials: creds);

return new JwtSecurityTokenHandler().WriteToken(token);
\end{lstlisting}
\end{english}

الشرح:
\begin{itemize}
\item \code{keyProvider.GetSigningKey()} — يجلب المفتاح \emph{الخاصّ}
  (\en{RSA private key})؛ في الإنتاج من \en{Azure Key Vault}، وفي التطوير مفتاح محلّي.
\item \code{SecurityAlgorithms.RsaSha256} — هذا هو \en{RS256} بالضبط.
\item \code{claims} — تحمل هويّة المستخدم وأدواره؛ لأنها موقّعة، أي تلاعب بها يُبطِل
  التوقيع فيُرفَض الرمز.
\end{itemize}

\begin{why}
لماذا \en{RS256}؟ لأن الخدمات التي \emph{تتحقّق} من الرمز تحتاج المفتاح
\emph{العامّ} فقط، والمفتاح العامّ لا يستطيع التوقيع؛ فحتى لو تسرّب لا يستطيع أحد
تزوير رموز. أمّا \en{HS256} فسرّه الواحد يوقّع \emph{ويتحقّق}، فمن يقرأه يقدر يزوّر.
وبفصل المفتاح خلف المنفذ \code{IJwtKeyProvider} صار \textbf{تدوير المفاتيح} عملية
تشغيلية بلا إعادة نشر للكود.
\end{why}

% ====================================================================
\subsection{الدخول الاجتماعي (SSO) والتحقّق بخطوتين (2FA)}

إلى جانب كلمة المرور، يدعم النظام الدخول عبر \en{Google} و\en{Microsoft} بمعيار
\en{OAuth 2.0} (تدفّق كود التفويض)، ويديره \code{SsoService}. الفكرة أن المستخدم
يوثّق نفسه عند \en{Google}، فيعود متصفّحه إلينا حاملًا \emph{كودًا} مؤقّتًا، ونحن
نبادل هذا الكود بمعلومات المستخدم في مكالمتين خلفيتين.

هنشوف قلب المبادلة من \code{SsoService.ExchangeCodeForTokenAsync} (مختصرًا):

\begin{english}
\begin{lstlisting}[language={[Sharp]C}]
var form = new Dictionary<string, string>
{
    ["code"]          = code,
    ["client_id"]     = provider.ClientId!,
    ["client_secret"] = provider.ClientSecret!,
    ["redirect_uri"]  = redirectUri,
    ["grant_type"]    = "authorization_code",
};

var response = await client.PostAsync(tokenEndpoint, content, ct);
// on success -> read the access_token, then GET the userinfo endpoint
// (email, name, picture) to build our SsoUserInfo
\end{lstlisting}
\end{english}

الشرح:
\begin{itemize}
\item \code{["code"]} — الكود المؤقّت العائد من المزوّد؛ يُبادَل مرّة واحدة فقط.
\item \en{client secret} — سرّ تطبيقنا لدى المزوّد، يأتي من الإعدادات لا من
  الكود، فلا يظهر أبدًا للمتصفّح.
\item بعد الحصول على \en{access token} نطلب نقطة \en{userinfo} لجلب البريد والاسم
  والصورة، ونبني منها \code{SsoUserInfo} لإنشاء الحساب أو ربطه.
\end{itemize}

\begin{secnote}
عنوان العودة (\en{redirect URI}) لا يُعكَس كما جاء من العميل؛ يفضّل النظام عنوانًا
ثابتًا من الإعدادات، أو يطابقه بقائمة سماح صريحة. هذا يسدّ ثغرة إعادة توجيه الكود إلى
موقع مهاجم — نقطة ضعف كلاسيكية في \en{OAuth}.
\end{secnote}

\begin{note}
للحسابات التي تريد حماية أقوى، يدعم النظام \textbf{التحقّق بخطوتين (\en{2FA})} عبر
\en{TOTP} — رمز لمرّة واحدة يتغيّر كل 30 ثانية من تطبيق مصادقة (مثل \en{Google
Authenticator}) — مبنيًّا على تدفّق \en{TwoFactor} في \en{ASP.NET Identity}. حتى لو
سُرِّبت كلمة المرور، يبقى الرمز الدوري حاجزًا ثانيًا.
\end{note}

% ====================================================================
\subsection{التفويض: الأدوار وفحوص الملكيّة}

بعد أن عرفنا \emph{مَن} المستخدم، يبدأ التفويض. يعمل على مستويين متكاملين:

\textbf{أولًا: الأدوار (\en{roles}).} لكل مستخدم دور أو أكثر، ويُحدَّد المسموح لكل دور:

\begin{center}
\begin{tabular}{L{3.4cm} L{11.3cm}}
\toprule
\textbf{\color{spPrimary}الدور} & \textbf{\color{spPrimary}صلاحيّته العامّة} \\
\midrule
\en{SuperAdmin} & أعلى صلاحية؛ إدارة المسؤولين والإعدادات الحسّاسة للمنصّة. \\
\en{Admin} & إدارة المحتوى والمستخدمين والمراجعات ومهامّ التشغيل اليومية. \\
\en{Student} & البحث عن المنح، التقديم، حجز الاستشارات، المشاركة في المجتمع. \\
\en{Consultant} & استقبال الحجوزات، عقد الجلسات، إدارة ملفّه المهني. \\
\en{Company} & نشر المنح ومراجعة الطلبات (مزوّد المنح). \\
\bottomrule
\end{tabular}
\end{center}

الأدوار تحمى نقاط النهاية على مستوى واسع (مثلًا: لوحة الإدارة لِـ\en{Admin} فقط).
لكن الدور وحده لا يكفي: كوْن المستخدم \en{Student} لا يعني أن له حقّ تعديل طلب
\emph{طالب آخر}.

\textbf{ثانيًا: فحوص الملكيّة (\en{ownership checks}).} داخل كل \en{handler} يتحقّق
النظام أن المورد المطلوب يخصّ صاحب الطلب فعلًا — بمقارنة مُعرِّف المستخدم الحالي
(من المنفذ \code{ICurrentUserService}) بمالك المورد — وإلّا يرفض العملية.

\begin{keyp}[نقطة مهمة: الدور طبقة أولى، والملكيّة طبقة ثانية]
الدور يجيب «هل يحقّ لهذا \emph{النوع} من المستخدمين استعمال هذه الوظيفة؟»، والملكيّة
تجيب «وهل يملك \emph{هذا} المستخدم \emph{هذا} المورد بالذات؟». الطبقتان معًا تمنعان
ثغرة الوصول المباشر لمورد غير مملوك (\en{IDOR}) — من أخطر أخطاء تطبيقات الويب.
\end{keyp}

% ====================================================================
\subsection{مخطّط تدفّق تسجيل الدخول}

نجمع كل ما سبق في تدفّق واحد: من وصول الطلب حتى إصدار الرمزين. اقرأ الجدول من أعلى
إلى أسفل:

\begin{center}
\begin{tabular}{L{1.1cm} L{13.4cm}}
\toprule
\textbf{\color{spPrimary}\#} & \textbf{\color{spPrimary}ما يحدث} \\
\midrule
1 & يصل طلب \en{POST} إلى نقطة الدخول ومعه البريد وكلمة المرور. \\
2 & يبحث \code{LoginCommandHandler} عن المستخدم بالبريد المُطبَّع؛ إن لم يوجد يُسجَّل فشل ويُرفَض برسالة عامّة. \\
3 & إن كان الحساب مقفولًا (\code{LockoutEnd} في المستقبل) يُرفَض فورًا. \\
4 & تُتحقّق كلمة المرور عبر \code{IPasswordHasher.Verify}؛ عند الفشل يُزاد العدّاد وقد يُقفَل الحساب (5 محاولات $\rightarrow$ 30 دقيقة). \\
5 & إن كان الحساب موقوفًا أو مُعطَّلًا يُرفَض الدخول رغم صحّة كلمة المرور. \\
6 & عند النجاح: تُصفَّر المحاولات، ويُحدَّث \code{LastLoginAt}، وتُسجَّل محاولة ناجحة. \\
7 & يجلب النظام أدوار المستخدم، ثم يستدعي \code{TokenService.IssueTokens}. \\
8 & يُنشأ \en{access JWT} موقّع بـ\en{RS256} (15 دقيقة) $+$ \en{refresh token} عشوائي. \\
9 & تُخزَّن \emph{تجزئة} رمز التحديث في \code{RefreshToken.TokenHash}، ويُعاد الزوج للعميل. \\
\bottomrule
\end{tabular}
\end{center}

\begin{qa}
\textbf{سؤال متوقّع:} «لو انتهى رمز الوصول والمستخدم لسّه شغّال، هيحصل إيه؟»\\
\textbf{الإجابة:} العميل يرسل رمز التحديث إلى نقطة التجديد، فيدوّره
\code{RotateRefreshTokenAsync}: يُبطِل القديم ويصدر زوجًا جديدًا، فتستمرّ الجلسة بلا
إدخال كلمة مرور. ولو ظهر رمز مُبطَل ثانيةً، تُلغى كل الجلسات كإجراء ضدّ السرقة.
\end{qa}

% ====================================================================
\begin{summary}[الخلاصة في سطور]
\begin{itemize}
\item \textbf{التوثيق} يجيب «مَن أنت؟»، و\textbf{التفويض} يجيب «تقدر تعمل إيه؟» —
  الأول عند الدخول، والثاني في كل طلب.
\item \code{ApplicationUser : IdentityUser<Guid>} يبني على \en{ASP.NET Core Identity}:
  كلمة مرور مُجزّأة بـ\en{PBKDF2}$+$مِلح، قفل بعد 5 محاولات، وتأكيد بريد.
\item النظام يُصدر \textbf{رمزين}: \en{access JWT} قصير (15 دقيقة) و\en{refresh token}
  يُدوَّر عند كل استخدام، ويُخزَّن \emph{مُجزَّأً} في \code{RefreshToken.TokenHash}.
\item نفس التجزئة تحمي \code{PasswordResetToken}: تسريب القاعدة لا يعطي رمزًا صالحًا.
\item التوقيع بـ\en{RS256} (غير متماثل) أأمن من \en{HS256}: المتحقّق يحتاج المفتاح
  العامّ فقط، والتدوير سهل عبر المنفذ \code{IJwtKeyProvider}.
\item \code{SsoService} يبادل أكواد \en{OAuth} مع \en{Google} و\en{Microsoft}،
  ويُتاح \en{2FA} عبر \en{TOTP}.
\item التفويض طبقتان: \textbf{الأدوار} (\en{SuperAdmin/Admin/Student/Consultant/Company})
  ثم \textbf{فحوص الملكيّة} داخل الـ\en{handlers} لسدّ ثغرات \en{IDOR}.
\end{itemize}
\end{summary}
